\chapter{Marco teórico}\label{chap:marco}
Este capítulo desarrolla el sustento científico que articula la óptica holográfica, la representación matemática de los hologramas y los modelos de aprendizaje automático empleados para su clasificación. Se pretende consolidar un hilo conductor que vaya desde la física de la formación del interferograma hasta los mecanismos de decisión, calibración y gobernanza que permiten utilizar el sistema en contextos clínicos reales.

\section{Formación del holograma y modelos de propagación}

En DLHM el sensor registra en una sola exposición la intensidad resultante de superponer una onda de referencia \(U_{\text{ref}}\) con la onda objeto \(U_{\text{obj}}\) modulada por la muestra. Asumiendo campo escalar, coherencia temporal suficiente y régimen paraxial, la intensidad en el plano del sensor \((x,y)\) es
\begin{equation}
I(x,y)=\bigl|U_{\text{ref}}(x,y)+U_{\text{obj}}(x,y)\bigr|^{2}
=|U_{\text{ref}}|^{2}+|U_{\text{obj}}|^{2}+2\,\Re\!\left\{U_{\text{ref}}\,U_{\text{obj}}^{*}\right\},
\label{eq:intensity}
\end{equation}
de donde se reconoce que el término de interferencia transporta la información de fase y amplitud de la muestra \cite{Amann2019}. En configuraciones de iluminación cuasi-plana, puede escribirse \(U_{\text{ref}}(x,y)=A_{\text{ref}}e^{j\phi_{\text{ref}}}\) con \(A_{\text{ref}}\) aproximadamente constante sobre el sensor; esta aproximación linealiza el contraste del interferograma y estabiliza el cociente señal-ruido (SNR) frente a fluctuaciones de \(U_{\text{obj}}\) \cite{Amann2019,OConnor2020}.

La propagación desde el plano del objeto \(z=0\) al plano del sensor \(z\) se modela, bajo hipótesis paraxial y banda limitada, por el integral de Fresnel:
\begin{equation}
U(x,y,z)=\frac{e^{jkz}}{j\lambda z}\iint U(x',y',0)\,
\exp\!\left\{\frac{jk}{2z}\!\left[(x-x')^{2}+(y-y')^{2}\right]\right\}\, \mathrm{d}x'\,\mathrm{d}y',
\label{eq:fresnel}
\end{equation}
con \(k=2\pi/\lambda\) \cite{Greenbaum2014}. Cuando la distancia muestra–sensor es pequeña respecto al tamaño del sensor o se requiere exactitud fuera de la paraxialidad, resulta más apropiada la formulación de Rayleigh–Sommerfeld,
\begin{equation}
U(x,y,z)=\frac{1}{j\lambda}\iint U(x',y',0)\,
\frac{e^{jkr}}{r}\,\cos\theta\;\mathrm{d}x'\,\mathrm{d}y',
\label{eq:RS}
\end{equation}
donde \(r=\sqrt{(x-x')^{2}+(y-y')^{2}+z^{2}}\) y \(\theta\) es el ángulo entre la normal al plano y el vector unión; el factor \(\exp(jkr)/r\) preserva la curvatura del frente de onda y permite cuantificar la sensibilidad del patrón a desalineaciones \cite{Greenbaum2014}. En prácticas computacionales, ambas expresiones se implementan mediante sus versiones convolutivas y FFT, cuidando la periodicidad implícita y el \emph{zero-padding} para evitar aliasing circular \cite{Greenbaum2014,OConnor2020}.

El muestreo discreto impone vínculos estrictos entre el \emph{pitch} del sensor \(\Delta\), la distancia \(z\) y el contenido angular del objeto. Si \(I\) se adquiere en una retícula \(\Delta\times\Delta\), la frecuencia espacial observable está acotada por \(f_{\max}=1/(2\Delta)\). Para un objeto con máximos ángulos de difracción \(\theta_{\max}\), la condición de Nyquist puede expresarse como
\begin{equation}
\frac{\lambda z}{\Delta^{2}} \gtrsim \pi^{-1}\quad\Longleftrightarrow\quad
z \gtrsim \frac{\pi\,\Delta^{2}}{\lambda},
\label{eq:nyquist}
\end{equation}
relación que fija cotas inferiores para \(z\) y evita que los órdenes de difracción se plieguen en el espectro del sensor \cite{Greenbaum2014}. La coherencia espacial se controla con diafragmas o guías de onda que introducen una función pupilar \(P(u,v)\) en el dominio de frecuencias; en presencia de esta apodización, el campo propagado se expresa como
\begin{equation}
U(x,y,z)=\mathcal{F}^{-1}\!\left\{\mathcal{F}\{U(x,y,0)\}\,H_{z}(u,v)\,P(u,v)\right\},
\label{eq:angular}
\end{equation}
donde \(H_{z}(u,v)\) es el filtro de transferencia de Fresnel en frecuencia. Esta forma pone de manifiesto cómo la óptica del sistema recorta componentes de alta frecuencia y justifica la normalización geométrica y radiométrica previa a la extracción de descriptores \cite{Greenbaum2014,OConnor2020}.

En configuraciones en línea (\emph{on-axis}) la medida \eqref{eq:intensity} contiene los términos \(|U_{\text{ref}}|^{2}\) y \(|U_{\text{obj}}|^{2}\) además del cruzado; la separación espectral de estos términos se logra explotando la curvatura de \(U_{\text{ref}}\) o con esquemas ligeramente \emph{off-axis} que desplazan el lóbulo de interferencia, facilitando el aislamiento del término portador sin reconstrucción explícita del campo complejo \cite{OConnor2020}. La cuantificación del SNR del interferograma puede aproximarse como
\begin{equation}
\mathrm{SNR}\approx
\frac{2A_{\text{ref}}A_{\text{obj}}}{\sqrt{\sigma_{\text{shot}}^{2}+\sigma_{\text{lectura}}^{2}+\sigma_{\text{cuant}}^{2}}},
\end{equation}
con varianzas asociadas a ruido de conteo, lectura y cuantización; esto motiva la elección de exposiciones que mantengan \(A_{\text{ref}}\) dominante pero sin saturación, optimizando el contraste de franjas \cite{Amann2019,OConnor2020}.

La estabilidad geométrica también es crítica: desalineaciones por \emph{tilt} o curvatura de frente se modelan como fases cuadráticas \(\exp\{j\pi (ax^{2}+by^{2}+cxy)\}\). Tales aberraciones se compensan estimando y retirando el plano/curvatura de fase en el dominio frecuencial, operación que mejora la estacionariedad de las franjas y reduce fugas espectrales \cite{OConnor2020}. En suma, las ecuaciones \eqref{eq:fresnel}–\eqref{eq:angular} y las restricciones de muestreo \eqref{eq:nyquist} delimitan la región de operación segura del montaje DLHM, aportando una base física para los pasos de control de calidad y normalización empleados aguas abajo \cite{Greenbaum2014,Amann2019,OConnor2020}.
\section{Representación matemática de texturas holográficas}

\begin{figure}[t]
    \centering
    \includegraphics[width=0.95\textwidth]{figuras/hologram_transforms.png}
    \caption{Transformaciones clave aplicadas sobre un holograma falciforme: interferograma original, magnitud logarítmica de la transformada de Fourier, mapa uniforme de LBP, respuesta de un filtro de Gabor, energía combinada de coeficientes wavelet y matriz de coocurrencia GLCM normalizada.}
    \label{fig:hologram-transforms}
\end{figure}

Un holograma en intensidad \(I:\Omega\subset\mathbb{Z}^{2}\to\mathbb{R}_{+}\) codifica modulaciones mixtas de amplitud y fase cuyo patrón de franjas exhibe periodicidades, anisotropías y discontinuidades finas. La estrategia adoptada es construir un vector de características \(\mathbf{x}\in\mathbb{R}^{d}\) que preserve esa física mediante descriptores locales y globales en dominios espacial, frecuencial y multiescala, evitando la reconstrucción del campo complejo. El flujo se resume en la Fig.~\ref{fig:hologram-transforms} y se fundamenta en operadores con propiedades de invariancia y estabilidad bien caracterizadas \cite{Amann2019,OConnor2020,BuitragoDuque2023}.

\paragraph{Preprocesamiento radiométrico y normalización geométrica.}
Sea \(\tilde I\) la versión corregida de \emph{flat-field} y \emph{dark-frame} de \(I\); se normaliza a rango unitario con
\[
I_{n}(x,y)=\frac{\tilde I(x,y)-\mathrm{P}_{1}\{\tilde I\}}
{\mathrm{P}_{99}\{\tilde I\}-\mathrm{P}_{1}\{\tilde I\}},
\]
donde \(\mathrm{P}_{q}\{\cdot\}\) denota el percentil \(q\). Para estabilizar el contraste local se emplea realce adaptativo tipo CLAHE sobre \(I_{n}\), lo que atenúa saturaciones y conserva bordes finos relevantes para la textura \cite{Amann2019,BuitragoDuque2023}. Todas las operaciones se realizan tras redimensionar a \(512\times512\) para fijar una malla de muestreo común y respetar las restricciones espectrales discutidas en la Sección previa \cite{Greenbaum2014}.

\paragraph{Patrones locales (LBP).}
La microtextura alrededor de \((x_{c},y_{c})\) se codifica con el operador
\begin{equation}
\operatorname{LBP}_{P,R}(x_{c},y_{c})=\sum_{p=0}^{P-1}
\mathbf{1}\!\left[g_{p}-g_{c}\ge 0\right]\,2^{p},
\label{eq:lbp}
\end{equation}
donde \(g_{c}=I_{n}(x_{c},y_{c})\) y \(\{g_{p}\}_{p=0}^{P-1}\) son intensidades en una circunferencia de radio \(R\) muestreada con interpolación bilineal. La estadística usada es el histograma normalizado \(h\in\Delta^{2^{P}-1}\). Para reducir la varianza y obtener invariancia a rotación, se adopta la versión uniforme \(u2\) y la variante \(\operatorname{LBP}_{ri}\) basada en rotaciones cíclicas que minimizan el patrón binario \cite{Dheepak2023}.

\paragraph{Coocurrencias de niveles de gris (GLCM).}
Sea \(Q\) la cuantización escalar de \(I_{n}\) a \(L\) niveles. Para un desplazamiento \(\boldsymbol{\Delta}=(\Delta_{x},\Delta_{y})\) y un ángulo \(\theta\), la matriz de coocurrencias \(P_{\boldsymbol{\Delta},\theta}\in[0,1]^{L\times L}\) se define por
\begin{equation}
P_{\boldsymbol{\Delta},\theta}(i,j)=\frac{1}{N}\sum_{(x,y)\in\Omega}
\mathbf{1}\!\left[Q(x,y)=i\right]\,
\mathbf{1}\!\left[Q(x+\Delta_{x},y+\Delta_{y})=j\right],
\label{eq:glcm}
\end{equation}
con normalización \(\sum_{i,j}P(i,j)=1\) \cite{Barburiceanu2021}. A partir de \eqref{eq:glcm} se obtienen medidas interpretables: contraste \(C=\sum_{ij}(i-j)^{2}P(i,j)\), homogeneidad \(H=\sum_{ij}P(i,j)/[1+(i-j)^{2}]\), energía \(E=\sum_{ij}P(i,j)^{2}\) y correlación \(\rho=\sum_{ij} ij\,P(i,j)-\mu_{x}\mu_{y}\), con medias \(\mu_{x}\), \(\mu_{y}\). Promediar sobre \(\theta\in\{0,\pi/4,\pi/2,3\pi/4\}\) atenúa efectos direccionales espurios y captura la granularidad interferométrica \cite{Barburiceanu2021}.

\paragraph{Contenido frecuencial y anisotropía angular.}
Sea \(\widehat I=\mathcal{F}\{I_{n}\}\) la transformada de Fourier bidimensional. La densidad espectral de potencia \(S(u,v)=|\widehat I(u,v)|^{2}\) permite definir un perfil radial mediante anillos concéntricos \(\mathcal{A}_{k}=\{(u,v): r_{k}\le\sqrt{u^{2}+v^{2}}<r_{k+1}\}\):
\begin{equation}
E_{k}=\iint_{\mathcal{A}_{k}} S(u,v)\, \mathrm{d}u\,\mathrm{d}v,
\label{eq:ringenergy}
\end{equation}
que resume el reparto de energía por número de franjas; hologramas con patrones más regulares tienden a concentrar energía en bajas frecuencias \cite{Amann2019,OConnor2020}. La anisotropía se cuantifica con el segundo momento angular del espectro,
\begin{equation}
\mathcal{A}=\frac{\lambda_{\max}-\lambda_{\min}}{\lambda_{\max}+\lambda_{\min}},\qquad
\{\lambda_{\min},\lambda_{\max}\}=\text{autovalores de } \mathbf{M}=\iint
\begin{bmatrix} u^{2} & uv \\ uv & v^{2}\end{bmatrix}
\frac{S(u,v)}{\iint S}\,\mathrm{d}u\,\mathrm{d}v,
\label{eq:anisotropy}
\end{equation}
donde \(\mathcal{A}\in[0,1)\) crece con la direccionalidad de las franjas \cite{Amann2019}.

\begin{figure}[H]
    \centering
    \includegraphics[width=0.95\textwidth]{figuras/descriptor_profiles.png}
    \caption{Comparativa de descriptores entre un holograma sano (fila superior) y uno falciforme (fila inferior). Cada columna exhibe el holograma original, el histograma LBP, la energía espectral por anillos de Fourier y la energía acumulada de los coeficientes wavelet.}
    \label{fig:descriptor-profiles}
\end{figure}

\paragraph{Filtros de Gabor.}
La respuesta orientada se modela con bancos \( \{g_{\theta,\lambda}\} \) definidos por
\begin{equation}
g_{\theta,\lambda}(x,y)=
\exp\!\left(-\frac{x'^{2}+\gamma^{2}y'^{2}}{2\sigma^{2}}\right)
\cos\!\left(\frac{2\pi x'}{\lambda}\right),\quad
\begin{cases}
x' = x\cos\theta + y\sin\theta,\\
y' = -x\sin\theta + y\cos\theta,
\end{cases}
\label{eq:gabor}
\end{equation}
y características basadas en momentos de \(|I_{n}*g_{\theta,\lambda}|\) (media, varianza y curtosis por par \((\theta,\lambda)\)) \cite{Alzubaidi2020}.

\paragraph{Energía multiescala por wavelets.}
La transformada wavelet discreta separa \(I_{n}\) en coeficientes de aproximación \(A^{(\ell)}\) y detalles \(D^{(\ell)}_{h},D^{(\ell)}_{v},D^{(\ell)}_{d}\) por nivel \(\ell\). La energía por escala se define como
\begin{equation}
W^{(\ell)}=\frac{1}{|D^{(\ell)}|}
\sum_{(x,y)}\!\left[D^{(\ell)}_{h}(x,y)^{2}+D^{(\ell)}_{v}(x,y)^{2}+D^{(\ell)}_{d}(x,y)^{2}\right],
\label{eq:waveenergy}
\end{equation}
y captura irregularidades finas y bordes a distintas resoluciones \cite{Alzubaidi2020}.

\paragraph{Morfometría en máscara binaria.}
Para complementar la textura se obtiene una máscara \(M\) por umbralización adaptativa y limpieza morfológica. A partir de \(M\) se calculan área \(A=\sum M\), perímetro \(P\) y casco convexo \(A_{\text{conv}}\). La circularidad \(C=4\pi A/P^{2}\), la solidez \(S=A/A_{\text{conv}}\) y la excentricidad \(e=\sqrt{1-\lambda_{2}/\lambda_{1}}\) (con \(\lambda_{1}\ge\lambda_{2}\) autovalores de la matriz de inercia) cuantifican deformaciones geométricas asociadas a falcización \cite{Alzubaidi2020}. Se añaden momentos de intensidad a primer–cuarto orden y entropía de Shannon sobre \(I_{n}\) para capturar dispersión radiométrica \cite{Singh2017}.

\paragraph{Estandarización y selección.}
Todas las características se estandarizan con estimadores robustos (mediana y MAD) y se filtran con pruebas de varianza casi nula; se aplica selección supervisada \(\mathrm{SelectKBest}\) con estadísticos monotónicos y validación anidada para evitar fuga de información \cite{Tsamardinos2022}. Este pipeline produce vectores \(\mathbf{x}\) estables frente a variaciones suaves de iluminación, traslación y rotación, preservando señales físicas relevantes del interferograma \cite{Amann2019,OConnor2020,BuitragoDuque2023}.
\section{Modelación estadística de los clasificadores}

La tarea de discriminación entre hologramas de eritrocitos sanos y falciformes se formula como un problema de inferencia probabilística en el que, dado un vector de características \(\mathbf{x}\in\mathbb{R}^{d}\), se desea estimar \(p(y=1\mid\mathbf{x})\) con \(y\in\{0,1\}\). El aprendizaje se plantea mediante la minimización del riesgo empírico regularizado, lo que implica resolver
\begin{equation}
\hat f \in \arg\min_{f\in\mathcal{H}} \ \frac{1}{n}\sum_{j=1}^{n}\ell\!\left(y_{j},f(\mathbf{x}_{j})\right)+\Omega(f),
\label{eq:erm}
\end{equation}
con una pérdida \(\ell\) apropiada para clasificación y un término \(\Omega\) que controla la complejidad para favorecer generalización, especialmente relevante en el régimen de pocos datos y alta dimensionalidad característico de DLHM \cite{Tsamardinos2022,Vabalas2019}. La estimación de hiperparámetros y la selección de transformaciones se realiza dentro de un esquema de validación cruzada anidada con particiones estratificadas, de forma que cualquier paso de selección o reescalado quede confinado al pliegue de entrenamiento y se evite fuga de información \cite{Tsamardinos2022,Vabalas2019}.

La regresión logística penalizada modela la probabilidad posterior mediante un mapeo sigmoidal de una combinación lineal de características, \(P(y=1\mid\mathbf{x})=\sigma(\beta_{0}+\sum_{i=1}^{d}\beta_{i}x_{i})\) con \(\sigma(z)=1/(1+e^{-z})\). La estimación se formula como una maximización de verosimilitud con penalización \(\ell_{1}\), \(\ell_{2}\) o combinada (elastic net), tal como
\begin{equation}
\mathcal{L}(\boldsymbol{\beta})=-\sum_{j=1}^{n}\!\left[y_{j}\log \sigma(z_{j})+(1-y_{j})\log\!\big(1-\sigma(z_{j})\big)\right]
+\lambda\!\left(\alpha\lVert\boldsymbol{\beta}\rVert_{1}+\frac{1-\alpha}{2}\lVert\boldsymbol{\beta}\rVert_{2}^{2}\right),
\label{eq:logit_elasticnet}
\end{equation}
donde \(\lambda>0\) controla la magnitud de la regularización y \(\alpha\in[0,1]\) pondera la esparsidad frente a la estabilidad numérica \cite{Pedregosa2011}. Esta formulación ofrece coeficientes interpretables en términos de log-odds y, con ponderaciones por clase cuando existe desbalance, permite ajustar sensibilidad y especificidad al criterio clínico.

Cuando la separación lineal en el espacio original no es adecuada, la máquina de soporte vectorial con núcleo gaussiano proyecta implícitamente las muestras a un espacio de Hilbert de dimensión infinita en el que la maximización del margen es efectiva. El problema dual
\begin{equation}
\max_{\boldsymbol{\alpha}}\ \sum_{j=1}^{n}\alpha_{j}-\frac{1}{2}\sum_{j,k=1}^{n}\alpha_{j}\alpha_{k}y_{j}y_{k}K(\mathbf{x}_{j},\mathbf{x}_{k}),
\quad \text{sujeto a } 0\le \alpha_{j}\le C,\ \sum_{j}\alpha_{j}y_{j}=0,
\label{eq:svm_dual}
\end{equation}
con \(K(\mathbf{x},\mathbf{x}')=\exp(-\gamma\lVert\mathbf{x}-\mathbf{x}'\rVert^{2})\), produce una función de decisión \(g(\mathbf{x})=\mathrm{sign}(\sum_{j\in SV}\alpha_{j}y_{j}K(\mathbf{x},\mathbf{x}_{j})+b)\) que controla el compromiso entre error de entrenamiento y complejidad mediante \(C\) y \(\gamma\) \cite{Pedregosa2011}. Para disponer de probabilidades calibradas, el puntaje de margen se transforma a escala probabilística con un ajuste logístico o isotónico en un conjunto de calibración independiente del ajuste del hiperparámetro \cite{Pedregosa2011}.

Los modelos basados en árboles exploran interacciones no lineales y umbrales en las variables sin requerir supuestos de linealidad. Un bosque aleatorio entrena múltiples árboles sobre remuestreos de las observaciones y subespacios aleatorios de características, y promedia sus salidas probabilísticas para reducir la varianza debida a la inestabilidad de cada árbol individual. La impureza (Gini o entropía) guía las particiones axis-aligned, y la agregación final hereda robustez ante \emph{outliers} y escalas heterogéneas \cite{Pedregosa2011}. En contraste, el gradiente por potenciación construye un modelo aditivo \(F_{m}(\mathbf{x})=F_{m-1}(\mathbf{x})+\nu\,h_{m}(\mathbf{x})\) donde \(h_{m}\) es un árbol poco profundo ajustado al gradiente negativo de la pérdida logística; la tasa de aprendizaje \(\nu\), el número de iteraciones y la profundidad controlan el equilibrio sesgo–varianza y actúan como regularización explícita en el espacio de funciones \cite{Friedman2001}. La combinación de ambos enfoques aporta capacidad de modelado no lineal complementaria a los clasificadores lineales.

Con el fin de estabilizar el desempeño en el régimen de pocos datos y reducir la varianza de estimación, las salidas probabilísticas de modelos heterogéneos se combinan en una mezcla convexa
\begin{equation}
p_{\mathrm{ens}}(y=1\mid\mathbf{x})=\sum_{k}w_{k}\,p_{k}(y=1\mid\mathbf{x}),\qquad w_{k}\ge 0,\ \sum_{k}w_{k}=1,
\label{eq:softvote}
\end{equation}
donde los pesos \(w_{k}\) se determinan en un pliegue de calibración para minimizar una métrica propia de probabilidad, como la pérdida de Brier, preservando la interpretabilidad agregada y amortiguando errores idiosincráticos de cada base \cite{Pedregosa2011}. La decisión final \(\hat y=\mathbf{1}[p_{\mathrm{ens}}\ge \tau]\) se ajusta a costes clínicos mediante la elección de un umbral \(\tau\) que optimiza, por ejemplo, el índice de Youden o que satisface una cota mínima de sensibilidad cuando el falso negativo es crítico.

La evaluación rigurosa de generalización se fundamenta en métricas umbral-independientes y umbral-dependientes. Se reporta el área bajo la curva ROC y, en presencia de desbalance, la curva de precisión–recobrado, junto con precisión balanceada, sensibilidad y especificidad puntuales. La incertidumbre de estas estimaciones se cuantifica con repetición de la validación cruzada estratificada y con intervalos de confianza por \emph{bootstrap}, siempre respetando la anidación: toda transformación de características, estandarización robusta y selección supervisada se ejecuta dentro de cada pliegue de entrenamiento para evitar sesgos optimistas \cite{Vabalas2019,Tsamardinos2022}. La calibración probabilística se verifica mediante diagramas de confiabilidad y pérdida de Brier, aplicando, cuando procede, calibración isotónica o Platt sobre un conjunto de calibración separado del ajuste de hiperparámetros \cite{Pedregosa2011}.

En conjunto, la articulación de un clasificador lineal penalizado, un margin-based kernel method y ensambles de árboles, todos integrados bajo una validación anidada y una capa explícita de calibración, proporciona un andamiaje estadístico coherente con las restricciones del problema: alta dimensionalidad relativa al tamaño muestral, potenciales colinealidades entre descriptores físicos y necesidad de interpretar decisiones a nivel de característica. Este marco da soporte a las afirmaciones de desempeño con garantías estadísticas explícitas y mantiene la trazabilidad necesaria para auditoría y gobernanza del modelo en contextos clínicos \cite{Tsamardinos2022,Vabalas2019,Pedregosa2011,Friedman2001}.


\section{Detección de anomalías, calibración e incertidumbre}

El funcionamiento seguro en un entorno clínico requiere distinguir entre entradas pertenecientes al dominio de entrenamiento y observaciones que se desvían de él, así como cuantificar y comunicar la incertidumbre asociada a cada decisión. La detección de anomalías se aborda mediante la distancia de Mahalanobis definida sobre el espacio de características estandarizadas. Sea \(\mathbf{x}\in\mathbb{R}^{d}\) el vector de características y sea \((\boldsymbol{\mu},\boldsymbol{\Sigma})\) la media y la covarianza robustas estimadas en el conjunto de entrenamiento; la estadística
\begin{equation}
D_{M}^{2}(\mathbf{x})=(\mathbf{x}-\boldsymbol{\mu})^{\top}\boldsymbol{\Sigma}^{-1}(\mathbf{x}-\boldsymbol{\mu})
\label{eq:mahala}
\end{equation}
resume la discrepancia cuadrática con respecto a la distribución empírica del dominio \cite{Mahalanobis1936}. Para atenuar la influencia de atípicos y colinealidades en \(\boldsymbol{\Sigma}\), la estimación se realiza con procedimientos robustos basados en el \emph{Minimum Covariance Determinant} (MCD), que maximizan la verosimilitud sobre subconjuntos de alta densidad y ofrecen consistencia bajo contaminación moderada \cite{Rousseeuw2019}. Bajo supuestos aproximadamente gaussianos en el espacio transformado, \(D_{M}^{2}\) se distribuye como \(\chi^{2}_{d}\), lo que permite fijar un umbral \(\tau_{\alpha}=\chi^{2}_{d,1-\alpha}\) para un nivel de significancia \(\alpha\) y derivar valores \(p\) mediante \(p(\mathbf{x})=1-F_{\chi^{2}_{d}}(D_{M}^{2}(\mathbf{x}))\). La descomposición por coordenadas
\begin{equation}
c_{i}(\mathbf{x})=\bigl(e_{i}^{\top}\boldsymbol{\Sigma}^{-1/2}(\mathbf{x}-\boldsymbol{\mu})\bigr)^{2},\qquad i=1,\dots,d,
\end{equation}
con \(\{e_{i}\}\) base canónica, ofrece una lectura de contribuciones que facilita interpretar qué componentes del descriptor impulsan la alerta y, por tanto, sustenta la revisión clínica de casos fuera de dominio \cite{Rousseeuw2019}. En la operación prospectiva, el umbral \(\tau_{\alpha}\) se valida por \emph{bootstrap} estratificado para mantener tasas de falsa alarma controladas en tamaños muestrales finitos y para documentar la estabilidad del detector a variaciones de calibración.

La incertidumbre asociada a la clasificación se cuantifica en dos niveles complementarios. En primer lugar, la calibración probabilística verifica la concordancia entre probabilidades pronosticadas y frecuencias observadas. Dada una función \(p(\mathbf{x})\in[0,1]\), la pérdida de Brier
\begin{equation}
\mathrm{BS}=\frac{1}{n}\sum_{j=1}^{n}\bigl(p(\mathbf{x}_{j})-y_{j}\bigr)^{2}
\end{equation}
proporciona una medida propia, y las curvas de confiabilidad contrastan visualmente la relación entre probabilidad estimada y tasa empírica de positivos en intervalos de puntuación. Cuando se detecta descalibración, se emplea ajuste isotónico o mapeo logístico sobre un conjunto de calibración separado del ajuste de hiperparámetros, manteniendo la independencia entre entrenamiento, calibración y prueba \cite{Pedregosa2011}. En segundo lugar, la predicción conforme introduce conjuntos de salida con garantía de cobertura finita-muestral bajo intercambiabilidad. Sea \(\alpha(\mathbf{x},y)\) una medida de no conformidad calculada a partir de las puntuaciones del modelo; en el esquema \emph{split conformal}, se calcula el cuantil \(q_{1-\epsilon}\) de las no conformidades en un conjunto de calibración y se define el conjunto de predicción
\begin{equation}
\Gamma_{\epsilon}(\mathbf{x})=\bigl\{y\in\{0,1\}:\alpha(\mathbf{x},y)\le q_{1-\epsilon}\bigr\},
\end{equation}
que garantiza cobertura marginal \(1-\epsilon\) sin supuestos paramétricos \cite{Montero2024}. En clasificación binaria, la elección \(\alpha(\mathbf{x},1)=1-p(\mathbf{x})\) y \(\alpha(\mathbf{x},0)=p(\mathbf{x})\) produce conjuntos \(\Gamma_{\epsilon}(\mathbf{x})\) de tamaño uno (decisión puntual) o dos (abstención informada), lo que resulta apropiado cuando los costos clínicos penalizan con severidad los falsos negativos. La combinación de calibración probabilística y predicción conforme habilita tanto la toma de decisiones costo–sensitivas mediante umbrales explícitos como la emisión de alertas de baja confianza documentadas y auditables \cite{Montero2024,Pedregosa2011}.

\section{Integración interdisciplinaria, ética y líneas de investigación}

El desempeño técnico del sistema depende de un acoplamiento explícito entre la modelación óptica, el diseño estadístico y la práctica clínica. La documentación de parámetros ópticos y condiciones de adquisición—distancia muestra–sensor, longitud de onda, paso del pixel y protocolos de control de iluminación—asegura la reproducibilidad de los interferogramas y sustenta las hipótesis de propagación y muestreo empleadas en el preprocesamiento. La construcción del vector de características y la selección de clasificadores se apoya en validación cruzada anidada y en procedimientos de calibración probabilística, con registro de todas las decisiones de ingeniería y de sus efectos sobre las métricas para mantener trazabilidad y permitir auditoría independiente \cite{Tsamardinos2022,Vabalas2019,Pedregosa2011}. La verificación clínica incorpora la revisión de casos límite, el estudio de errores por subgrupos y la evaluación del impacto de umbrales costo–sensitivos en escenarios de cribado, privilegiando configuraciones con sensibilidad mínima garantizada cuando el falso negativo comporta riesgo sanitario significativo.

El cumplimiento normativo y la gobernanza del modelo se alinean con el reglamento de propiedad intelectual de la Universidad EAFIT, que exige preservar la trazabilidad de datos y modelos, delimitar claramente los derechos de uso y asegurar supervisión humana en decisiones asistenciales \cite{EAFIT2018}. En paralelo, las recomendaciones contemporáneas sobre inteligencia artificial responsable enfatizan la transparencia metodológica, la explicabilidad de las decisiones y los mecanismos de detección de entradas fuera del dominio, todos ellos incorporados mediante reportes reproducibles, análisis de contribuciones en la distancia de Mahalanobis y predicción conforme con cobertura controlada \cite{Montero2024}. Esta arquitectura sugiere líneas de investigación continuas: la caracterización óptica detallada del montaje facilitaría la generación de datos sintéticos controlados para estudios de sensibilidad; los modelos híbridos que integren descriptores físicos con representaciones aprendidas explicables podrían mejorar la sensibilidad a patrones sutiles manteniendo interpretabilidad; la validación multicéntrica permitiría cuantificar robustez ante cambios de dominio y sostener generalización externa en poblaciones diversas.

La consolidación de estos componentes articula un marco que vincula la física de la formación del holograma con la representación estadística y la inferencia probabilística, y que incorpora salvaguardas de detección de dominio y comunicación de incertidumbre. Este marco no sólo proporciona garantías cuantitativas sobre el desempeño esperado, sino que también define prácticas de documentación, supervisión humana y auditoría que son coherentes con la integración del sistema en flujos de trabajo clínicos reales \cite{EAFIT2018,Montero2024}.
